%==============================================================
\section{Θεωρητικό Υπόβαθρο}
%==============================================================

Η παρούσα ενότητα παρουσιάζει τις θεμελιώδεις θεωρητικές έννοιες και το μαθηματικό υπόβαθρο που υποστηρίζουν τη σχεδίαση συστημάτων Ανάκτησης-Επαυξημένης Παραγωγής (Retrieval-Augmented Generation - RAG) σε πολύστροφα (multi-turn) διαλογικά περιβάλλοντα. Συγκεκριμένα, αναλύεται η στοχαστική φύση των Μεγάλων Γλωσσικών Μοντέλων, το υβριδικό παραδειγματικό πλαίσιο της αρχιτεκτονικής RAG, οι σύγχρονες μέθοδοι νευρωνικής ανάκτησης πληροφορίας, οι τεχνικές αναδιατύπωσης ερωτημάτων μέσω συλλογιστικής, καθώς και οι προκλήσεις της συνομιλιακής αναζήτησης και οι μεθοδολογίες πολυδιάστατης αξιολόγησής τους.

%--------------------------------------------------------------
\subsection{Μεγάλα Γλωσσικά Μοντέλα}
%--------------------------------------------------------------

Τα Μεγάλα Γλωσσικά Μοντέλα (ΜΓΜ - Large Language Models) αποτελούν αυτο-παλινδρομικά (autoregressive) στοχαστικά μοντέλα γλώσσας, βασισμένα στην αρχιτεκτονική Transformer και ειδικότερα στον μηχανισμό της αυτο-προσοχής (self-attention). Τα μοντέλα αυτά εκπαιδεύονται μέσω μεγιστοποίησης της λογαριθμικής πιθανότητας του επόμενου συμβόλου (next-token prediction) πάνω σε εκτεταμένα σώματα κειμένου, μια διαδικασία που τυπικά μοντελοποιείται ως:

\begin{equation}
P(X) = \prod_{t=1}^{T} P(x_t \mid x_1, x_2, \dots, x_{t-1}) = \prod_{t=1}^{T} P(x_t \mid x_{<t})
\end{equation}

Η εκθετική κλιμάκωση (scaling laws) του αριθμού των παραμέτρων, του όγκου των δεδομένων εκπαίδευσης και της υπολογιστικής ισχύος οδηγεί στην εμφάνιση αναδυόμενων ικανοτήτων (emergent abilities). Οι σημαντικότερες από αυτές περιλαμβάνουν:
\begin{itemize}
    \item \textbf{Μάθηση εντός πλαισίου (In-Context Learning - ICL):} Η ικανότητα εκτέλεσης νέων εργασιών κατά τη φάση της εξαγωγής συμπερασμάτων (inference) αποκλειστικά μέσω της παροχής παραδειγμάτων στην προτροπή (prompt), χωρίς τροποποίηση των βαρών του δικτύου.
    \item \textbf{Συμμόρφωση με οδηγίες (Instruction Following):} Η ευθυγράμμιση του μοντέλου με τις προθέσεις του χρήστη μέσω εκπαίδευσης εποπτευόμενης λεπτομερούς ρύθμισης (Supervised Fine-Tuning - SFT) και ενισχυτικής μάθησης από ανθρώπινη ανάδραση (RLHF).
    \item \textbf{Πολυβηματικός Συλλογισμός (Chain-of-Thought Reasoning):} Η ικανότητα αποδόμησης σύνθετων προβλημάτων σε διαδοχικά, λογικά ενδιάμεσα βήματα συλλογιστικής πριν από την παραγωγή της τελικής απόκρισης. Στο Κεφάλαιο 5 αξιοποιούμε άμεσα αυτή την αρχή στη σχεδίαση της στρατηγικής Chain-of-Thought αναδιατύπωσης ερωτήματος (Ενότητα 5.2.1) και στην πολυσταδιακή διαδικασία παραγωγής απάντησης (Ενότητα 5.3).
\end{itemize}

Παρά την υψηλή τους αναπαραστατική ικανότητα, τα ΜΓΜ παρουσιάζουν δομικούς περιορισμούς:
\begin{itemize}
    \item \textbf{Στατικότητα γνώσης (Knowledge Cutoff):} Η ενσωματωμένη γνώση παραμένει χρονικά και πραγματολογικά περιορισμένη στο σημείο ολοκλήρωσης της προεκπαίδευσης.
    \item \textbf{Έλλειψη εγγενούς ανιχνευσιμότητας (Traceability):} Τα μοντέλα λειτουργούν ως "μαύρα κουτιά" (black boxes), αδυνατώντας να παρέχουν ρητές παραπομπές ή αποδείξεις για την ορθότητα των ισχυρισμών τους.
    \item \textbf{Πραγματολογικές Ψευδαισθήσεις (Hallucinations):} Η τάση παραγωγής κειμένου που είναι συντακτικά και υφολογικά άρτιο, αλλά πραγματολογικά εσφαλμένο ή ανυπόστατο, λόγω της φύσης τους ως προσεγγιστές κατανομών πιθανότητας.
\end{itemize}

Οι περιορισμοί αυτοί επιβάλλουν τη μετάβαση από αμιγώς παραμετρικά συστήματα σε υβριδικές αρχιτεκτονικές που ενσωματώνουν εξωτερικές, δυναμικές πηγές γνώσης.

%--------------------------------------------------------------
\subsection{Μέθοδοι Ανάκτησης Πληροφορίας}
%--------------------------------------------------------------

Οι σύγχρονες αρχιτεκτονικές ανάκτησης πληροφορίας κατηγοριοποιούνται με βάση τον τρόπο αναπαράστασης του κειμένου:

\begin{itemize}
    \item \textbf{Λεξιλογικοί Ανακτητές (Sparse Retrievers):} Βασίζονται στην ακριβή αντιστοίχιση όρων (exact keyword matching) χρησιμοποιώντας στατιστικές συχνότητας εμφάνισης. Ο αλγόριθμος BM25 αποτελεί το τυπικό βιομηχανικό πρότυπο, βαθμολογώντας τη συνάφεια ενός εγγράφου $d$ με βάση τον όρο $t$ ως συνάρτηση της συχνότητας του όρου εντός του εγγράφου (term frequency) και της αντίστροφης συχνότητας εμφάνισης στο σώμα κειμένου (inverse document frequency). Παρά την υψηλή ακρίβεια σε τεχνικούς όρους και σειριακούς αριθμούς, το BM25 περιορίζεται θεμελιωδώς από το πρόβλημα της ασυμφωνίας λεξιλογίου (vocabulary mismatch): αν το ερώτημα χρησιμοποιεί παραφράσεις ή συνώνυμα που δεν συμπίπτουν λεξιλογικά με το κείμενο-στόχο, η τομή $q \cap d$ μηδενίζεται.
    \item \textbf{Πυκνοί Ανακτητές (Dense Retrievers):} Χρησιμοποιούν βαθιά νευρωνικά δίκτυα (π.χ. αρχιτεκτονικές Dual-Encoder) εκπαιδευμένα μέσω αντιπαραθετικής μάθησης (contrastive learning) για τη χαρτογράφηση ερωτημάτων και εγγράφων σε έναν κοινό, χαμηλής διάστασης διανυσματικό χώρο $\mathbb{R}^d$. Η βαθμολογία συνάφειας υπολογίζεται μέσω του εσωτερικού γινομένου ή της συνημίτονου ομοιότητας των αντίστοιχων embeddings: $S(q, d) = \langle \mathbf{E}_Q(q), \mathbf{E}_D(d) \rangle$. Παρότι αποτελεσματικοί σε αφηρημένη σημασιολογική ταύτιση και διαγλωσσικές παραφράσεις, εισάγουν συχνά σημασιολογικό θόρυβο σε πολύστροφα περιβάλλοντα, υπεργενικεύοντας τεχνικά tokens.
    \item \textbf{Μαθητοί Αραιοί Ανακτητές (Learned Sparse Retrievers):} Αποτελούν μια σύγχρονη υβριδική προσέγγιση (π.χ. SPLADE, ELSER) που ενοποιεί τα πλεονεκτήματα λεξιλογικής ακρίβειας και σημασιολογικής γενίκευσης. Χρησιμοποιούν ένα γλωσσικό μοντέλο για να προβλέψουν και να επεκτείνουν το λεξιλόγιο ενός εγγράφου (term expansion) στον αρχικό χώρο των high-dimensional tokens, επιλύοντας το πρόβλημα της συνωνυμίας χωρίς να χάνουν την ακρίβεια της αραιής αναπαράστασης. Η βαθμολογία ορίζεται ως:
    \begin{equation}
    S(q, d) = \sum_{t \in q \cap d} w_q(t) \cdot w_d(t)
    \end{equation}
    όπου τα βάρη $w(t)$ είναι μαθημένες μη αρνητικές τιμές που αντανακλούν τη σημασιολογική βαρύτητα του όρου εντός του συγκεκριμένου corpus.
\end{itemize}

Ο Πίνακας~\ref{tab:retrieval-comparison} συνοψίζει τα τρία αυτά παραδείγματα ως προς τα κρίσιμα χαρακτηριστικά τους για το πολύστροφο σενάριο. Η παρούσα εργασία βασίζεται σε έναν corpus-aligned learned sparse ανακτητή (\textbf{ELSER v1}), ο οποίος, όπως αναλύεται διεξοδικά στην Ενότητα 1.5, υπερτερεί εμπειρικά τόσο του BM25 όσο και των ισχυρότερων εμπορικών dense ανακτητών (π.χ. Cohere Embed v4) στο benchmark MTRAG, με το πλεονέκτημα να αποδίδεται εν μέρει στη μείωση του vocabulary gap χωρίς τον σημασιολογικό θόρυβο των πυκνών embeddings (με την επιφύλαξη της μεροληψίας των κρίσεων συνάφειας, Ενότητα 1.5).

\begin{table}[h]
\centering
\small
\begin{tabular}{lccc}
\toprule
\textbf{Χαρακτηριστικό} & \textbf{Sparse (BM25)} & \textbf{Dense} & \textbf{Learned Sparse (ELSER v1)} \\
\midrule
Αναπαράσταση & Λεξιλογική, αραιή & Νευρωνική, πυκνή ($\mathbb{R}^d$) & Λεξιλογική, αραιή με term expansion \\
Χειρισμός συνωνύμων & Όχι & Ναι & Ναι (μέσω επέκτασης όρων) \\
Ακρίβεια σε IDs/κωδικούς & Υψηλή & Χαμηλή & Υψηλή \\
Ερμηνευσιμότητα & Πλήρης & Περιορισμένη & Υψηλή \\
Ευθυγράμμιση με corpus & — & Γενικού σκοπού & Corpus-aligned \\
Επίδοση στο MTRAG (R@5) & 0.153 & 0.397 (καλύτερο dense) & \textbf{0.483} \\
\bottomrule
\end{tabular}
\caption{Συγκριτική επισκόπηση παραδειγμάτων ανάκτησης. Οι τιμές Recall@5 στη στήλη MTRAG προέρχονται από την αξιολόγηση χωρίς αναδιατύπωση ερωτήματος (777 ερωτήματα development set) και αιτιολογούν την επιλογή του ELSER v1 ως βάσης ανάκτησης του προτεινόμενου συστήματος.}
\label{tab:retrieval-comparison}
\end{table}

%--------------------------------------------------------------
\subsection{Αναδιατύπωση Ερωτημάτων και Σύντηξη}
%--------------------------------------------------------------

Σε δυναμικά διαλογικά περιβάλλοντα, τα ερωτήματα των χρηστών είναι συχνά ελλιπή ή εξαρτώνται άμεσα από το ιστορικό. Η Αναδιατύπωση Ερωτημάτων (Query Reformulation) μέσω ΜΓΜ μετασχηματίζει ένα εξαρτημένο ερώτημα σε μια αυτοτελή (context-independent) μορφή. Οι κυριότερες στρατηγικές περιλαμβάνουν:
\begin{itemize}
    \item \textbf{HyDE (Hypothetical Document Embeddings):} Παραγωγή μιας υποθετικής απάντησης (zero-shot) από το ΜΓΜ, η οποία στη συνέχεια χρησιμοποιείται ως διευρυμένο ερώτημα για την ανάκτηση, στοχεύοντας στην ευθυγράμμιση στον χώρο των απαντήσεων.
    \item \textbf{Chain-of-Thought (CoT) Reformulation:} Παραγωγή ενδιάμεσων βημάτων ανάλυσης της πληροφοριακής ανάγκης πριν την τελική εξαγωγή των όρων αναζήτησης.
    \item \textbf{Keyword Extraction \& Anchor Targeting:} Εστίαση στην εξαγωγή ονοματοποιημένων οντοτήτων και λέξεων-κλειδιών που λειτουργούν ως άγκυρες (anchors) στο συγκεκριμένο σώμα κειμένου.
\end{itemize}

Όταν εκδίδονται πολλαπλές, συμπληρωματικές αναδιατυπώσεις ερωτήματος, απαιτείται ένας μηχανισμός σύντηξης για τη συγχώνευση των επιμέρους κατατάξεων. Η Σύντηξη Αντίστροφης Κατάταξης (Reciprocal Rank Fusion - RRF) αποτελεί μια ανθεκτική μέθοδο (unsupervised rank aggregation) που ορίζεται ως:

\begin{equation}
\text{RRF\_Score}(d \in \mathcal{D}) = \sum_{m \in \mathcal{M}} \frac{1}{k + r_m(d)}
\end{equation}

όπου $\mathcal{M}$ είναι το σύνολο των παραχθέντων λιστών κατάταξης, $r_m(d)$ είναι η θέση (rank) του εγγράφου $d$ στη λίστα $m$, και $k$ είναι μια σταθερά εξομάλυνσης (smoothing constant, τυπικά $k=60$). Η RRF εξασφαλίζει ότι έγγραφα που κατατάσσονται ψηλά σε πολλαπλές στρατηγικές ευνοούνται, χωρίς να επηρεάζονται από τις αποκλίσεις στις απόλυτες τιμές των σκορ των διαφορετικών μοντέλων. Ωστόσο, η επίπεδη (flat) εκδοχή της RRF αντιμετωπίζει όλες τις στρατηγικές αναδιατύπωσης ως ισοδύναμες, κάτι που μπορεί να είναι προβληματικό όταν ορισμένες (π.χ. HyDE, CoT) αναμένεται να είναι λιγότερο σταθερές ανά στροφή από άλλες (π.χ. Minimal, Corpus-Specific). Μια ιεραρχική παραλλαγή της RRF που αντιμετωπίζει αυτή την ανομοιογένεια παρουσιάζεται στην Ενότητα 1.4.2.

%--------------------------------------------------------------
\subsection{Πολύστροφη Συνομιλιακή RAG}
%--------------------------------------------------------------

Σε αντίθεση με το μονοστροφικό RAG, η Πολύστροφη Συνομιλιακή RAG (Multi-Turn Conversational RAG) εισάγει τη διάσταση του χρόνου και της διαλογικής εξάρτησης. Το κύριο πρόβλημα εντοπίζεται στη διατήρηση της \textit{σταθερότητας ανάκτησης συμφραζομένων} (context-aware retrieval stability). Καθώς ο διάλογος εξελίσσεται, τα ερωτήματα εμφανίζουν:
\begin{itemize}
    \item \textbf{Επίλυση Αναφορών (Coreference Resolution):} Χρήση αντωνυμιών ή δεικτικών εκφράσεων που αναφέρονται σε οντότητες προηγούμενων στροφών.
    \item \textbf{Ελλειπτικότητα (Ellipsis):} Παράλειψη ολόκληρων συντακτικών δομών ή υποκειμένων που θεωρούνται γνωστά από τα συμφραζόμενα.
    \item \textbf{Μετατόπιση Θεματικού Πλαισίου (Topic Switching):} Απότομες αλλαγές στη ροή της συζήτησης.
\end{itemize}

Η μη επιτυχής διαχείριση αυτών των φαινομένων οδηγεί σε \textit{σωρευτική διάδοση σφαλμάτων} (error cascading): ένα αποτυχημένο βήμα ανάκτησης στη στροφή $t$ εισάγει θόρυβο στη γεννήτρια, η οποία παράγει μια εσφαλμένη απάντηση, η οποία με τη σειρά της μολύνει το ιστορικό του διαλόγου, προκαλώντας ακόμα μεγαλύτερη απόκλιση (context drift) στη στροφή $t+1$. Ο μηχανισμός αυτός προϋποθέτει ότι το ιστορικό περιέχει τις προηγούμενες απαντήσεις του ίδιου του συστήματος, όπως σε έναν πραγματικό βοηθό· το MTRAGEval, αντίθετα, παρέχει σε κάθε στροφή το ιστορικό αναφοράς, οπότε η αξιολόγησή μας δεν τον ενεργοποιεί. Στο MTRAG benchmark παρατηρούμε πάντως απότομη πτώση της απόδοσης ανάκτησης μετά την πρώτη στροφή (Ενότητα 1.5.2).

%--------------------------------------------------------------
\subsection{Αξιολόγηση Συστημάτων RAG}
%--------------------------------------------------------------

Η αξιολόγηση σε πολύστροφα περιβάλλοντα RAG επιβάλλει τη χρήση ενός πολυδιάστατου πλαισίου μετρικών που εξετάζει ανεξάρτητα την ανάκτηση και την παραγωγή. Οι βασικοί άξονες περιλαμβάνουν:
\begin{itemize}
    \item \textbf{Μετρικές Ανάκτησης εγγράφων:} Το Recall@k (ικανότητα εντοπισμού όλων των σχετικών εγγράφων) και το Normalized Discounted Cumulative Gain (nDCG@k), το οποίο λαμβάνει υπόψη τη σχετική θέση και τη διαβαθμισμένη συνάφεια των ανακτηθέντων αποσπασμάτων.
    \item \textbf{Μετρικές Πιστότητας και Παραγωγής:} Το αναφορικά αλγοριθμικό σκορ $RB_{alg}$ (σύνθεση BERT-Recall, BERT-K-Precision και ROUGE-L έναντι της ανθρώπινα επιμελημένης αναφοράς) και το χωρίς αναφορά μετρικό πιστότητας $RL_F$ (πλαίσιο RAGAS), που αποτιμά αν οι ισχυρισμοί της απάντησης θεμελιώνονται στα ανακτηθέντα αποσπάσματα.
    \item \textbf{Αξιολόγηση μέσω ΜΓΜ (LLM-as-a-Judge):} Το αναφορικό σκορ $RB_{llm}$, ένας διάμεσος-συγκεντρωμένος αξιολογητής πολλαπλών ΜΓΜ που εκτιμά Πιστότητα (Faithfulness), Καταλληλότητα (Appropriateness) και Πληρότητα (Completeness) της απάντησης.
\end{itemize}

Στο πλαίσιο του MTRAGEval (SemEval-2026 Task 8), η end-to-end απόδοση του συστήματος συντίθεται μέσω του Αρμονικού Μέσου (Harmonic Mean - HM) των επιμέρους αυτών διαστάσεων:

\begin{equation}
\text{HM} = \frac{3}{\frac{1}{\text{RB}_{alg}} + \frac{1}{\text{RL}_F} + \frac{1}{\text{RB}_{llm}}}
\end{equation}

Η χρήση του αρμονικού μέσου τιμωρεί αυστηρά συστήματα που παρουσιάζουν ασυμμετρία στην απόδοσή τους (π.χ. υψηλό F1 αλλά χαμηλό Faithfulness λόγω ψευδαισθήσεων), επιβάλλοντας την ισορροπημένη βελτιστοποίηση όλων των συνιστωσών. Σημειώνεται επιπλέον ότι η αξιολόγηση εξαρτάται κρίσιμα από έναν αυτοματοποιημένο ταξινομητή «Δεν Γνωρίζω» (IDK classifier): σωστές αρνήσεις απάντησης σε μη απαντήσιμες στροφές λαμβάνουν τέλειο σκορ (1.0) σε όλους τους άξονες, ενώ εσφαλμένες αρνήσεις ή εσφαλμένες αποδοχές μηδενίζονται — καθιστώντας την πύλη απαντησιμότητας υψηλού ρίσκου συνιστώσα του συστήματος, όπως αναλύεται στην Ενότητα 1.5.4.

%--------------------------------------------------------------
\subsection{Σύνοψη Θεωρητικού Υποβάθρου}
%--------------------------------------------------------------

Συνοψίζοντας, η διασταύρωση των Μεγάλων Γλωσσικών Μοντέλων με τις τεχνικές Νευρωνικής Ανάκτησης Πληροφορίας δημιουργεί ισχυρά υβριδικά συστήματα RAG, η σταθερότητα των οποίων δοκιμάζεται σε πολύστροφα σενάρια διαλόγου. Ο περιορισμός της αστάθειας της ανάκτησης μεταξύ στροφών ---και, σε έναν πραγματικό βοηθό, της σωρευτικής διάδοσης σφαλμάτων--- απαιτεί τον αυστηρό έλεγχο της ποιότητας των ερωτημάτων μέσω πολλαπλών στρατηγικών αναδιατύπωσης, τη χρήση εύρωστων αλγορίθμων σύντηξης κατατάξεων και την εισαγωγή πυλών απαντησιμότητας. Η παρούσα εργασία βασίζεται πάνω σε αυτές τις αρχές για τη διασφάλιση της σταθερότητας της ανάκτησης και της πιστότητας των παραγόμενων αποκρίσεων.